\section{Consequences and further questions}\label{sec:discussion}
The two main results answer different questions. The corner expansion
explains the size of the norm and the origin of its lower-order terms.
The maximizer theorem says why that particular corner value is the norm.
Only eventual corner maximality is needed to transfer an asymptotic
expansion; the finite completion supplies the stronger, all-degree
geometric statement.

The estimate in Theorem~\ref{thm:main-expansion} implies
\[
 N\left\{\Lambda_n-\frac4{\pi^2}(\log N+d)^2-\kappa\right\}
 \longrightarrow -1.
\]
Thus the reciprocal correction is distinguishable from an unspecified
vanishing remainder. The proof does not give a sharp numerical constant
in its final $O(\log N/N^2)$ error, nor determine the next coefficient.
The analytic onset $320$ is sufficient for the localization argument;
we do not claim that it is the smallest possible onset.

There are natural questions beyond the present result. One can seek a
further asymptotic term, or a fully analytic maximizer proof covering the
small degrees without a finite computation. It would also be useful to
understand which parts of the envelope-minus-defect comparison extend
to other Lissajous interpolation sets. The Morrow--Patterson family,
for example, requires a separate analysis of its Lebesgue constant. The present corner formula does not by itself establish such
an extension.
